% ECONOMICS_PROBLEM: {"id": "Q58-500", "title": "Carbon-price responsiveness", "jel_code": "Q58", "source_id": "OP-0435"}
% !TeX program = xelatex
\documentclass[11pt,a4paper]{article}
\usepackage[margin=23mm]{geometry}
\usepackage{fontspec}
\setmainfont{Latin Modern Roman}
\usepackage{amsmath,amssymb,mathtools}
\usepackage{xcolor,enumitem,booktabs,longtable,array,xurl,hyperref}
\definecolor{muted}{HTML}{53636C}
\hypersetup{colorlinks=true,urlcolor=blue,linkcolor=blue}
\setlength{\parindent}{0pt}
\setlength{\parskip}{5pt}
\newcommand{\R}{\mathbb R}
\newcommand{\E}{\mathbb E}
\newcommand{\N}{\mathbb N}
\newcommand{\ind}{\mathbf 1}
\newcommand{\Var}{\operatorname{Var}}
\newcommand{\Cov}{\operatorname{Cov}}
\newcommand{\supp}{\operatorname{supp}}
\DeclareMathOperator*{\argmin}{arg\,min}
\DeclareMathOperator*{\argmax}{arg\,max}
\newcommand{\entrymeta}[3]{{\sffamily\small\color{muted}\textbf{\texttt{#1}}\quad|\quad #2\par #3\par}}
\newcommand{\Problemref}[1]{\texttt{#1}}
\newcommand{\JELref}[2]{\texttt{#2} (JEL \texttt{#1})}
\begin{document}
\section*{Carbon-price responsiveness}
\textbf{Problem ID:} \texttt{Q58-500}\quad \textbf{JEL:} \texttt{Q58}\par
% BEGIN ECONOMICS STATEMENT
\label{problem:OP-0435}
{\sffamily\small\color{muted}Original subject group: Climate, environment and energy\par}
\label{definitions:problem:30007717}
\textbf{Audit classification:} Published research agenda. Checked source identity and appropriate agenda/background support; expanded intervention, units, population, definedness and causal restrictions. Exact targets and variants are editorial.
\subsection*{Definitions and precise research target}
\textit{Editorial formalization.} Consider manufacturing plants covered by a prespecified carbon-pricing scheme during its first ten operating years. Let \(p>0\) be the real carbon price, \(E_i(p,s)\geq0\) plant emissions, and \(s\) the complete schedule of other policies, free allowances, and monitoring rules. Define aggregate baseline-cohort emissions \(M(p,s)=E[\sum_i E_i(p,s)]\), retaining closed plants with zero emissions and separately accounting for entry. Assume \(M(p_0,s)>0\) and differentiability near a specified policy price \(p_0\). The target elasticity is
\[\varepsilon(p_0,s)=\left.\frac{p}{M(p,s)}\frac{\partial M(p,s)}{\partial p}\right|_{p=p_0}.\]
The question is to identify this causal local response, and its two-year and ten-year versions, rather than infer it from comparisons of policy adoption alone. Specify a credible source of price variation and exclusion restrictions separating carbon prices from coincident regulation or economic shocks. Account for anticipation, allowance banking, output relocation, and emissions displaced to uncovered plants. If actual price variation cannot separately identify extensive and intensive adjustment, report their joint response or honest bounds.

For each horizon \(h\in\{2,10\}\), define \(E_{ih}(p,s)\) as cumulative tonnes of a stated emissions measure during years 1 through \(h\), under a constant real price \(p\) and a fixed complete policy schedule \(s\). Set \(M_h(p,s)=E\sum_{i\in I_0}E_{ih}(p,s)\), with the fixed baseline plant cohort \(I_0\); price-induced closures contribute zero. Define \(\varepsilon_h=(p_0/M_h(p_0,s))\partial_pM_h(p_0,s)\), assuming a positive denominator and differentiability. Entry emissions and displaced foreign or uncovered emissions are separate estimands. A tax variation and a cap-and-trade allowance-price variation are different interventions: under an exactly binding fixed cap, covered total emissions may be fixed even when plant behavior changes. All expectations use a stated baseline population and probability law. Identification means every admissible data-generating model with the same observed-data law yields the same displayed target; otherwise report the identified set. Exact equations, horizons and assumptions are editorial, rather than source-given canonical conjectures.
\subsection*{Variants and assumption changes}
\begin{enumerate}
\item \textbf{Editorial leakage variant.} Include entrants, uncovered and foreign emissions in \(M_h^{\rm global}(p,s)\); estimate its elasticity separately from covered baseline plants. Signs need not coincide.
\item \textbf{Editorial policy-path variant.} Replace constant \(p\) by an announced future vector \(\mathbf p\) and estimate \(\partial\log M_h/\partial\log p_t\), allowing expectations and allowance banking.
\end{enumerate}
\subsection*{References and source audit}
\textbf{1.} Explicit gap in price-level elasticity evidence, distinct from introduction-of-policy effects. Locator: Nature Communications 15, 4147; abstract p.1 and introduction p.2. Full text or primary publication page inspected. URL: \url{https://thedocs.worldbank.org/en/doc/83e2d33f048039a281426fa3bb150895-0050022025/original/Systematic-review-and-meta-analysis-of-ex-post-evaluations-on-the-effectiveness-of-carbon-pricing.pdf}.
\begin{enumerate}\raggedright
\item\hypertarget{definitions:source:30007717:1}{} Niklas Döbbeling-Hildebrandt; Klaas Miersch; Tarun M. Khanna; Marion Bachelet; Stephan B. Bruns; Max Callaghan; Ottmar Edenhofer; Christian Flachsland; Piers M. Forster; Matthias Kalkuhl; Nicolas Koch; William F. Lamb; Nils Ohlendorf; Jan Christoph Steckel; Jan C. Minx. 2024. \emph{Systematic review and meta-analysis of ex-post evaluations on the effectiveness of carbon pricing}. Nature Communications 15, 4147; abstract p.1 and introduction p.2.
\url{https://thedocs.worldbank.org/en/doc/83e2d33f048039a281426fa3bb150895-0050022025/original/Systematic-review-and-meta-analysis-of-ex-post-evaluations-on-the-effectiveness-of-carbon-pricing.pdf}.
DOI: \url{https://doi.org/10.1038/s41467-024-48512-w}.
\end{enumerate}

\subsection*{Common conventions: Source mathematical conventions}

These conventions apply to each entry unless explicitly replaced.

\textbf{Models and identification.} A primitive model \(m\) specifies all probability laws, feasible choices, information, equilibrium and measurement rules used by its target. The admissible class \(\mathcal M\) is fixed independently of outcomes. If \(P_O(m)\) is its observable probability law and \(\tau(m)\) the target, its identified set at observable law \(P\) is
\[
 \mathcal I_\tau(P)=\{\tau(m):m\in\mathcal M,\ P_O(m)=P\}.
\]
Point identification means that this set is a singleton. An empty set signals incompatibility of data and assumptions. Coordinate bounds need not describe the joint set. Bounds are sharp only if every retained target is attainable or arbitrarily approachable by compatible models. Finite bounds require explicit restrictions on unobserved quantities; unrestricted confounding, hidden wealth, extrapolation or arbitrary equilibrium selection can yield uninformative sets.

\textbf{Causal interventions.} A potential outcome \(Y_i(p)\) is the outcome under a complete policy regime \(p\), including timing, financing, information and market spillovers. The target population and units are fixed before treatment. With interference, use the complete assignment vector or a justified exposure mapping; individual no-interference notation alone cannot identify an economy-wide intervention. A difference of mean potential outcomes does not require the cross-world joint distribution. Conditioning on post-treatment migration, survival, enrollment or employment changes the estimand and needs separate justification.

\textbf{Statistical statements.} An estimator, confidence set, forecast or test specifies a sampling law, model class and loss/coverage criterion. Uniform \((1-\alpha)\)-coverage means \(\inf_{m\in\mathcal M}\Pr_m\{\tau(m)\in C_n\}\geq1-\alpha\), or an explicitly asymptotic version. The whole parameter space trivially has coverage, so informativeness requires a stated width, risk or power objective. A forecasting improvement is relative to a named benchmark, information set and evaluation population.

\textbf{Equilibrium and optimization.} An equilibrium comprises feasible plans, optimality, beliefs where needed, budget constraints and all market/resource clearing. The primitives or a stated benchmark define each correspondence \(\mathcal E(p;m)\). Nonemptiness is assumed or proved, never inferred just from compact policy sets. A selected-equilibrium target uses a declared selection rule; a robust target quantifies over all permitted equilibria. Use suprema or infima unless attainment is established. An \(\varepsilon\)-optimal policy is feasible and within \(\varepsilon\) of the finite optimum value. Report an empty feasible set separately.

\textbf{Welfare and accounting.} Normative utility, interpersonal normalization, social weights, numeraire, discounting and the use of public receipts are primitives. Behavioral utility need not coincide with the welfare criterion. Household welfare includes ownership income and rents once; adding profits again to compensating variation generally double-counts them. A monetary equivalent is defined by an explicit transfer and price convention, with monotonicity and range conditions. Average effects, mechanism contributions, transfers and welfare losses are different targets. Infinite sums require convergence and dynamic feasible plans require terminal or no-Ponzi conditions.

\textbf{Source versions.} A working-paper date, revision date and journal publication year are distinguished. A DOI for a journal version is not automatically the DOI of an earlier manuscript. Locators refer to the stated version and printed pages where specified. A metadata-only check does not verify a claim about a full-text conclusion. Every entry's source subsection narrows the provenance claim accordingly.
% END ECONOMICS STATEMENT
\end{document}
